\documentclass[11pt,a4paper]{article}

\usepackage[utf8]{inputenc}
\usepackage[T1]{fontenc}
\usepackage{amsmath}
\usepackage{amssymb}
\usepackage{amsthm}
\usepackage{geometry}

\geometry{margin=2.5cm}

\title{
Vehicle Machine Language (VML)
\\
Draft White Paper Version 0.1
}

\author{
Working Draft
}

\begin{document}

\begin{abstract}

This document proposes Vehicle Machine Language (VML), a machine-native
behavior representation for autonomous vehicles and advanced driver
assistance systems.

VML is motivated by the observation that natural language is highly
effective for human communication but poorly suited for direct machine
execution.
Unlike natural language, VML provides explicit semantics,
deterministic interpretation and formal verification capabilities.

The purpose of VML is not to replace human requirements writing.
Instead, VML serves as the target representation into which human
requirements are translated.

Future systems may employ Large Language Models (LLMs) or
Vision-Language Models (VLMs) during requirements engineering to
perform this translation.
However, no runtime component of the vehicle is expected to consume
natural language directly.

This document defines the first conceptual version of VML and provides
sufficient detail to enable an initial implementation.

\end{abstract}

\maketitle
\begin{center}
\textbf{Origin:} https://twg-fas.org
\end{center}

\section{Introduction}

Modern autonomous systems process increasingly complex information.

A typical driving decision may depend on:

\begin{itemize}
\item observed objects,
\item traffic regulations,
\item temporal constraints,
\item mission goals,
\item safety requirements,
\item uncertainty estimates.
\end{itemize}

Most existing systems represent these concepts using separate
subsystems.

VML attempts to provide a common formal representation for all
behavior-relevant information.

The long-term vision is:

\[
Natural\ Language
\rightarrow
VML
\rightarrow
Vehicle\ Behavior
\]

where natural language exists exclusively during engineering and VML
becomes the native runtime language.

\section{Design Principles}

VML is built around six principles.

\subsection{P1: No Runtime Natural Language}

No runtime decision shall depend on natural-language interpretation.

\[
Runtime \cap NaturalLanguage = \emptyset
\]

\subsection{P2: Explicit Semantics}

Every VML statement shall have one unique interpretation.

\subsection{P3: Composability}

Complex behavior shall be constructed from simpler behavior.

\subsection{P4: Verifiability}

Every VML specification shall be suitable for formal analysis.

\subsection{P5: Human Readability}

While machine-oriented, VML should remain understandable to engineers.

\subsection{P6: Technology Independence}

VML shall not depend on any specific perception system,
planning algorithm or neural network architecture.

\section{VML Overview}

VML consists of six layers.

\[
VML
=
L
\oplus
T
\oplus
M
\oplus
P
\oplus
G
\oplus
D
\]

where

\begin{align}
L & : Logical Layer \\
T & : Temporal Layer \\
M & : Metric Layer \\
P & : Probabilistic Layer \\
G & : Goal Layer \\
D & : Deontic Layer
\end{align}

Each layer answers a distinct question.

\begin{center}
\begin{tabular}{ll}
Layer & Question \\
\hline
Logic & What exists? \\
Temporal & When must it happen? \\
Metric & How much? \\
Probability & How certain are we? \\
Goal & What do we want? \\
Deontic & What is allowed? \\
\end{tabular}
\end{center}

\section{Logical Layer}

The logical layer defines facts.

Example:

\begin{verbatim}
Pedestrian(P1)
Vehicle(V4)
Lane(L2)
\end{verbatim}

Interpretation:

\begin{itemize}
\item object P1 is a pedestrian
\item object V4 is a vehicle
\item object L2 is a lane
\end{itemize}

Relationships are represented as predicates.

\begin{verbatim}
InLane(V4,L2)
AheadOf(P1,Ego)
\end{verbatim}

The logical layer provides the vocabulary used by the remaining
layers.

\section{Temporal Layer}

The temporal layer defines behavior over time.

VML adopts four basic temporal operators.

\subsection{Always}

\[
\Box \varphi
\]

Meaning:

\begin{quote}
``\(\varphi\) must always hold.''
\end{quote}

Example:

\[
\Box(\neg Collision)
\]

\subsection{Eventually}

\[
\Diamond \varphi
\]

Meaning:

\begin{quote}
``\(\varphi\) must eventually become true.''
\end{quote}

Example:

\[
\Diamond ParkingCompleted
\]

\subsection{Next}

\[
\bigcirc \varphi
\]

Meaning:

\begin{quote}
``\(\varphi\) must hold in the next state.''
\end{quote}

\subsection{Until}

\[
\varphi U \psi
\]

Meaning:

\begin{quote}
``Maintain \(\varphi\) until \(\psi\) becomes true.''
\end{quote}

Example:

\[
Braking
U
PedestrianClear
\]

\section{Metric Layer}

Driving requires continuous values.

Examples include distance, velocity and acceleration.

VML introduces built-in arithmetic expressions.

Examples:

\[
Distance(Ego,P1) > 2m
\]

\[
Velocity(Ego) < 30km/h
\]

\[
LateralAcceleration(Ego) < 0.3g
\]

Metric expressions may be combined with temporal operators.

Example:

\[
\Box
(
Distance(Ego,P1) > 2m
)
\]

\section{Probabilistic Layer}

Perception is uncertain.

Instead of representing facts as strictly true or false, VML
associates confidence values.

Example:

\[
Probability(Pedestrian(P1))
=
0.95
\]

VML introduces confidence thresholds.

Example:

\[
Require_{0.99}
(
SafeDistance
)
\]

Interpretation:

Safety shall hold with at least 99\% confidence.

\section{Goal Layer}

Goals describe desired future states.

Example:

\begin{verbatim}
Goal ParkingCompleted
\end{verbatim}

or

\begin{verbatim}
Goal At(Ego,ChargingStation)
\end{verbatim}

Goals differ from constraints.

Constraint:

\[
\Box(\neg Collision)
\]

Goal:

\[
\Diamond ParkingCompleted
\]

A valid plan must satisfy all constraints while attempting to achieve
all goals.

\section{Deontic Layer}

The deontic layer represents obligations, permissions and
prohibitions.

These concepts map naturally to traffic regulations.

Obligation:

\[
O(YieldToPedestrian)
\]

Permission:

\[
P(ChangeLane)
\]

Prohibition:

\[
F(CrossSolidLane)
\]

Traffic regulations should preferably be represented using deontic
constructs rather than procedural code.

\section{Example}

Natural language requirement:

\begin{quote}
The vehicle shall stop for pedestrians and remain stopped until the
crossing is clear.
\end{quote}

VML representation:

\[
O(YieldToPedestrian)
\]

\[
\Box
(
PedestrianCrossing
\rightarrow
(
Braking
U
CrossingClear
)
)
\]

This representation is substantially more explicit and directly
machine-processable.

\section{Example Runtime Architecture}

A reference implementation may follow the architecture:

\[
Sensors
\rightarrow
Perception
\rightarrow
VML\ State
\rightarrow
Planner
\rightarrow
Controller
\]

The perception system updates predicates.

The planner searches for actions satisfying all constraints.

The controller executes the selected action.

\section{VML Core Profile}

The first implementable VML profile is intentionally limited.

Required features:

\begin{itemize}
\item predicates,
\item temporal operators,
\item metric expressions,
\item goals,
\item obligations.
\end{itemize}

Optional future features:

\begin{itemize}
\item probabilistic reasoning,
\item utility functions,
\item learning components,
\item ontology inheritance,
\item multi-agent reasoning.
\end{itemize}

This subset is sufficient for constructing an initial parser,
runtime monitor and planner.

\section{Conclusion}

VML is proposed as a machine-native language for autonomous systems.

The central philosophy is simple:

\[
NaturalLanguage
\neq
RuntimeLanguage
\]

Human requirements are authored in natural language.

They are translated into VML.

Runtime systems operate exclusively on VML.

In this manner natural language assumes the role of source code,
while VML becomes the executable behavioral representation used within
the vehicle.

Future versions of this document will extend the language,
formalize syntax and define reference execution semantics.

\end{document}
